% 分野別類題: 複素指数関数の積分・sinc関数（2019 問4 型）
% オイラーの公式を用いた複素指数積分、有限区間の積分から sinc 関数を導く問題 3 問。
% コンパイル: TEXINPUTS="../../../00_common:$TEXINPUTS" latexmk -xelatex problems.tex
\documentclass[a4paper,11pt]{article}
\usepackage{preamble}

\begin{document}

\kamokumei{類題演習 --- Complex Exponential Integrals and the Sinc Function}

\shijibun{以下の問題では、オイラーの公式 $e^{i\theta} = \cos\theta + i\sin\theta$ を用いて複素指数関数の積分を計算し、実部・虚部を比較することで実積分の値を求めよ。必要であれば $\mathrm{sinc}(x) = \dfrac{\sin x}{x}$（$\mathrm{sinc}(0)=1$）の記法を用いてよい。}

\shomon{問1.}{実数 $k \neq 0$ と正の定数 $a$ に対して、次の積分を計算せよ。}
\[
\int_{-a}^{a} e^{-ikx}\,dx
\]
さらに、その結果が $2a\,\mathrm{sinc}(ka)$ と表せることを示せ。

\shomon{問2.}{実数 $a, b$（$a \neq 0$）に対して、次の積分を求めよ。}
\[
\int e^{ax}\cos(bx)\,dx
\]

\shomon{問3.}{高さ $1$、幅 $2a$ の矩形パルス関数
\[
f(x) =
\begin{cases}
1 & (|x| \le a) \\
0 & (|x| > a)
\end{cases}
\]
のフーリエ変換
\[
F(k) = \int_{-\infty}^{\infty} f(x)\,e^{-ikx}\,dx
\]
を求め、$F(k)$ が $\mathrm{sinc}$ 関数を用いて表せることを示せ。また $\displaystyle\lim_{k \to 0} F(k)$ を求めよ。}

\end{document}
